\documentclass[]{report}

\usepackage{import}

\import{../}{settings.tex}

\title{Cours d'anneaux et corps pour l'agrégation}

\author{Cours de Diego IZQUIERDO, transcrit par SADR TAHOURI Luca}
\date{}

\begin{document}

\maketitle
\tableofcontents

\chapter{Anneaux et algèbres}

\section{Définitions et exemples}

\begin{df}{Anneaux, anneaux commutatifs}{}
    Un anneau est un ensemble $A$ muni de deux lois de compositions interne:
    \[\fundfto{+}{A\times A}{A}{(a,b)}{a+b}\qquad \fundfto{\times }{A\times A}{A}{(a,b)}{a+b}\]
    appelées respectivement addition et multiplication, telles que:
    \begin{itemize}
        \item Le couple $(A,+)$ est un groupe abbélien, dont l'élément neutre est généralement noté $0_A$
        \item La multiplication est associative:
        \[\forall (a,b,c)\in A^3,(a\times b)\times c=a\times (b\times c)\]
        \item La multiplication possède un élément neutre, généralement noté $1_A$:
        \[\forall a\in A, 1_A\times a=a\times 1_A=a\]
        \item La multiplication est distributive par rapport à l'addition, c'est-à-dire:
        \[\forall (a, b, c) \in A^3, \quad (a + b) \times c = a \times b + a \times c, \qquad a \times (b + c) = a \times b + a \times c.\]
    \end{itemize}
    Un anneau est dit commutatif si la multiplication est commutative:
    \[\forall (a,b)\in A^2, a\times b=b\times a\]
\end{df}

\begin{ex}{}{}
    \begin{enumerate}
        \item L'anneau nul $A = \{0_A\}$ est un anneau commutatif. Un anneau $A$ est nul si, et seulement si, $0_A = 1_A$.
        \item Les ensembles $\Z$, $\Z/n\Z$, $\Q$ et $\C$ munis de l'addition et de la multiplication usuelles sont des anneaux commutatifs.
        \item Le produit direct $\prod_{i \in I} A_i$ d'une famille d'anneaux $(A_i, +, \times)_{i \in I}$ est un anneau (pour les lois d'addition et de multiplication terme à terme).
        \item Si $A$ est un anneau commutatif et $n$ un entier strictement positif, l'ensemble $\Mr_n(A)$ des matrices $n \times n$ à coefficients dans $A$ muni de l'addition et de la multiplication usuelles
        \[(a_{ij})_{1 \le i,j \le n} + (b_{ij})_{1 \le i,j \le n} := (a_{ij} + b_{ij})_{1 \le i,j \le n}\]
        \[(a_{ij})_{1 \le i,j \le n} (b_{ij})_{1 \le i,j \le n} := \left( \sum_{k=1}^n a_{ik} b_{kj} \right)_{1 \le i,j \le n}\]
        est un anneau. L'anneau $\Mr_n(A)$ n'est commutatif que si $A$ est l'anneau nul ou si $n = 1$.
    \end{enumerate}
\end{ex}

\begin{df}{Diviseurs de zéro, élément inversible}{}
    Un élément non-nul $a\in A$ est dit diviseur de zéro s'il existe $b\in A\wt \{0_A\}$ tel que $ab=0_A$ ou $ba=0_A$. L'anneau $A$ est dit intègre s'il est commutatif et non nul et s'il n'a pas de diviseur de zéro.\\
    Un élément $a\in A$ est dit inversible s'il existe $b\in A$ tel que $ab=ba=1_A$. L'ensemble des éléments inversibles de $A$ est un groupe, généralement noté $A^\times$. L'anneau $A$ est dit être un corps s'il est commutatif non nul et si tout élément non nul de $A$ est inversible.
\end{df}

\begin{df}{}{}
    
\end{df}

\section{Anneaux et polynômes}

\subsection{Polynômes en une indéterminée}

\subsection{Polynômes en plusieurs indéterminées}

\section{Algèbres}

\section{Morphismes d'anneaux et d'algèbres}

\section{Sous-anneaux et idéaux}

\section{Divisibilité et arithmétique dans un anneau intègre}

\section{Idéaux premiers et maximaux}

\section{Anneaux euclidiens}

\section{Anneaux factoriels}

\section{Retour aux anneaux de polynômes}

\subsection{Racines d'un polynôme à une variable}

\subsection{Factoralité des anneaux de polynômes}

\subsection{Polynômes symétriques}

\section{Fractions rationnelles}

\end{document}